% ============================================================
%  NWU Faculty of Engineering – Final Year Project Report
%  LaTeX Template  |  Based on official report format
%
%  HOW TO CUSTOMISE:
%  1. Fill in the \newcommand fields in the "PROJECT INFO" block.
%  2. Replace placeholder text inside each section.
%  3. Add/remove sections as required by your project.
%  4. Place your university logo as "nwu_logo.png" in the
%     same folder, or change \logofile to point elsewhere.
%  5. Compile twice with pdflatex (or use latexmk -pdf).
% ============================================================

% ── Document class ──────────────────────────────────────────
\documentclass[12pt, a4paper]{report}

% ── Packages ────────────────────────────────────────────────
\usepackage[a4paper,
            top=25mm, bottom=25mm,
            left=25mm, right=25mm]{geometry}

\usepackage{graphicx}          % figures
\usepackage{amsmath, amssymb}  % maths
\usepackage{booktabs}          % professional tables
\usepackage{array}             % extended column types
\usepackage{longtable}         % multi-page tables
\usepackage{caption}           % caption formatting
\usepackage{subcaption}        % sub-figures
\usepackage{float}             % [H] placement
\usepackage{enumitem}          % list control
\usepackage{hyperref}          % clickable links / bookmarks
\usepackage{xcolor}            % colour support
\usepackage{titlesec}          % section heading styles
\usepackage{titletoc}          % TOC styles
\usepackage{fancyhdr}          % headers/footers
\usepackage{setspace}          % line spacing
\usepackage{microtype}         % better justification
\usepackage{parskip}           % space between paragraphs
\usepackage{listings}          % code listings
\usepackage{siunitx}           % SI units
\usepackage{cite}              % numeric citations [1]
\usepackage{url}               % URLs in references
\usepackage{pdfpages}          % include external PDFs
\usepackage{appendix}          % appendix environment
\usepackage{tocloft}           % fine-tune TOC/LOF/LOT

% ── Colour palette (match NWU purple) ───────────────────────
\definecolor{nwupurple}{RGB}{100, 28, 128}   % NWU brand purple
\definecolor{darkgray}{RGB}{64, 64, 64}

% ── Hyperref setup ──────────────────────────────────────────
\hypersetup{
    colorlinks  = true,
    linkcolor   = black,
    citecolor   = black,
    urlcolor    = nwupurple,
    pdftitle    = {Final Year Project Report},
    pdfauthor   = {Student Name},
    bookmarksnumbered = true,
}

% ── Line spacing ────────────────────────────────────────────
\setstretch{1.0}   % single spacing (body justified, as in template)

% ── Paragraph spacing (no indent, small gap) ────────────────
\setlength{\parindent}{0pt}
\setlength{\parskip}{6pt plus 2pt minus 1pt}

% ── Section heading styles ──────────────────────────────────
%    Chapter  → "1  Introduction"  (large, bold)
%    Section  → "1.1  Background"  (medium bold)
%    Subsec   → "1.1.1  Detail"    (bold)
\titleformat{\chapter}[hang]
  {\normalfont\large\bfseries}
  {\thechapter}{1em}{}
\titlespacing*{\chapter}{0pt}{12pt}{6pt}

\titleformat{\section}[hang]
  {\normalfont\normalsize\bfseries}
  {\thesection}{1em}{}
\titlespacing*{\section}{0pt}{10pt}{4pt}

\titleformat{\subsection}[hang]
  {\normalfont\normalsize\bfseries}
  {\thesubsection}{1em}{}
\titlespacing*{\subsection}{0pt}{8pt}{2pt}

\titleformat{\subsubsection}[hang]
  {\normalfont\normalsize\bfseries}
  {\thesubsubsection}{1em}{}
\titlespacing*{\subsubsection}{0pt}{6pt}{2pt}

% Remove "Chapter N" prefix from chapter titles
\renewcommand{\chaptername}{}
\renewcommand{\thechapter}{\arabic{chapter}}

% ── Table / figure caption style ────────────────────────────
%    Tables:  "Table 1. Caption"  above the table, centred, italic
%    Figures: "Figure 1. Caption" below the figure, centred, italic
\captionsetup[table]{
    labelsep   = period,
    labelfont  = {normalsize},
    textfont   = {it, normalsize},
    position   = above,
    justification = centering,
    singlelinecheck = false,
}
\captionsetup[figure]{
    labelsep   = period,
    labelfont  = {normalsize},
    textfont   = {it, normalsize},
    position   = below,
    justification = centering,
}

% ── Header / footer ─────────────────────────────────────────
\pagestyle{fancy}
\fancyhf{}
\renewcommand{\headrulewidth}{0pt}   % no header rule (matches template)
\fancyfoot[R]{\thepage}              % page number bottom-right
\fancypagestyle{plain}{%             % also affects chapter first pages
  \fancyhf{}
  \fancyfoot[R]{\thepage}
  \renewcommand{\headrulewidth}{0pt}
}

% ── TOC depth ───────────────────────────────────────────────
\setcounter{tocdepth}{3}
\setcounter{secnumdepth}{4}

% ── Listings (code) style ───────────────────────────────────
\lstset{
    basicstyle   = \ttfamily\small,
    breaklines   = true,
    frame        = single,
    numbers      = left,
    numberstyle  = \tiny\color{darkgray},
    keywordstyle = \color{nwupurple}\bfseries,
    commentstyle = \color{darkgray}\itshape,
    captionpos   = b,
    tabsize      = 4,
}

% ────────────────────────────────────────────────────────────
%    PROJECT INFO – edit these fields only  
% ────────────────────────────────────────────────────────────
\newcommand{\reporttitle}{Your Project Title Here}
\newcommand{\studentname}{Mr/Ms Firstname Lastname}
\newcommand{\studentnumber}{12345678}
\newcommand{\supervisorname}{Dr/Prof Supervisor Name}
\newcommand{\degree}{Bachelor of Engineering}
\newcommand{\discipline}{Mechatronic Engineering}   % e.g. Electrical / Civil / Mechanical
\newcommand{\campus}{Potchefstroom Campus}
\newcommand{\submissiondate}{November 2025}
\newcommand{\logofile}{nwu_logo.png}   % path to university logo file
\newcommand{\turnitinscore}{X\%}       % plagiarism score
% ────────────────────────────────────────────────────────────

% ============================================================
\begin{document}
% ============================================================

% ── COVER PAGE ──────────────────────────────────────────────
\begin{titlepage}
  \thispagestyle{empty}

  % Purple diagonal banner (top-left corner decoration)
  \begin{picture}(0,0)
    \put(-28, -80){%
      \color{nwupurple}%
      \rule{210mm}{70mm}%
    }
  \end{picture}

  % Logo in the top-left on the purple banner
  \vspace*{-12mm}
  \hspace*{-5mm}%
  %\includegraphics[height=22mm]{\logofile}

  \vspace{30mm}

  % Faculty heading (purple)
  \begin{center}
    {\Large\bfseries\color{nwupurple} FACULTY OF ENGINEERING}

    \vspace{18mm}

    % Report title
    {\large \reporttitle}

    \vspace{28mm}

    % Author block
    {\bfseries By:}\\[4pt]
    {\bfseries \studentname}\\[2pt]
    Student number: \studentnumber

    \vspace{18mm}

    {\bfseries Supervisor:}\\[4pt]
    {\bfseries \supervisorname}

    \vspace{18mm}

    Submitted in pursuit of the degree

    \vspace{8pt}
    {\bfseries \MakeUppercase{\degree}}\\[4pt]
    In\\[4pt]
    {\bfseries \MakeUppercase{\discipline}}

    \vspace{14mm}

    {\bfseries North-West University \campus}

    \vfill
    \thepage   % page 1 (printed bottom-right by pagestyle, mimic cover)
  \end{center}
\end{titlepage}

% ── FRONT MATTER (roman page numbers) ───────────────────────
\pagenumbering{roman}

% ─── Abstract ───────────────────────────────────────────────
\chapter*{Abstract}
\addcontentsline{toc}{chapter}{Abstract}

This report presents the design, development, and testing of \ldots{}
[Replace this with a concise 200–350 word summary covering: the
problem addressed, design approach, key results, and main conclusions.]

% ─── Keywords ───────────────────────────────────────────────
\chapter*{Keywords}
\addcontentsline{toc}{chapter}{Keywords}

Keyword One, Keyword Two, Keyword Three, Keyword Four, Keyword Five

% ─── Table of Contents ──────────────────────────────────────
\tableofcontents

% ─── List of Figures ────────────────────────────────────────
\listoffigures
\addcontentsline{toc}{chapter}{List of Figures}

% ─── List of Tables ─────────────────────────────────────────
\listoftables
\addcontentsline{toc}{chapter}{List of Tables}

% ─── Declaration of Originality ─────────────────────────────
\chapter*{Declaration of Originality}
\addcontentsline{toc}{chapter}{Declaration of Originality}

I, \studentname, with student number \studentnumber{} declare the following:

This report is my own original work. I further declare that:

\begin{enumerate}
    \item No part of it has been copied from another person;
    \item I did not work with another person on this project and report;
    \item I acknowledged all consulted sources in the text and submitted
          a bibliography;
    \item I have indicated all sections where I have used generative
          AI-enabled technology;
    \item Parts without references are my own ideas, arguments, and
          conclusions.
\end{enumerate}

The Turnitin plagiarism test resulted in a score of \turnitinscore.

\vspace{20mm}
\noindent\rule{60mm}{0.4pt}\\[2pt]
Signature\\[12pt]
\submissiondate

% ── MAIN BODY (arabic page numbers) ─────────────────────────
\cleardoublepage
\pagenumbering{arabic}

% ============================================================
%  CHAPTER 1 – INTRODUCTION
% ============================================================
\chapter{Introduction}
\label{ch:introduction}

This report outlines the design, development, and testing of \ldots{}
[Provide context for the project, explain the motivation, and
describe the overall structure of the report.]

The report is structured as follows: Chapter~\ref{ch:problem} defines
the problem and requirements. Chapter~\ref{ch:literature} presents
the literature study. Chapters~\ref{ch:concept}–\ref{ch:detailed}
cover the concept, preliminary, and detailed design phases.
Chapter~\ref{ch:implementation} describes the implementation.
Chapter~\ref{ch:testing} presents the testing and results.
Chapter~\ref{ch:conclusion} provides the conclusions.

% ============================================================
%  CHAPTER 2 – PROBLEM ANALYSIS
% ============================================================
\chapter{Problem Analysis}
\label{ch:problem}

\section{Background}
\label{sec:background}

[Describe the domain context and motivation for the project. Explain
why the problem exists and why it is important to solve it.]

\section{Problem Statement}
\label{sec:problem-statement}

[Clearly state the engineering problem to be solved. Be specific and
concise – one or two paragraphs.]

\section{Project Objectives}
\label{sec:objectives}

The objectives of this project are:

\begin{itemize}
    \item Objective 1: \ldots
    \item Objective 2: \ldots
    \item Objective 3: \ldots
\end{itemize}

\section{Requirements}
\label{sec:requirements}

\subsection{Functional Requirements}
\label{subsec:func-req}

\begin{itemize}
    \item \textbf{REQ\_01\_FUNC}: [Requirement description.]
    \item \textbf{REQ\_02\_FUNC}: [Requirement description.]
    \item \textbf{REQ\_03\_FUNC}: [Requirement description.]
\end{itemize}

\subsection{Performance Requirements}
\label{subsec:perf-req}

\begin{itemize}
    \item \textbf{REQ\_04\_PERF}: [Requirement description.]
    \item \textbf{REQ\_05\_PERF}: [Requirement description.]
\end{itemize}

\subsection{Physical Requirements}
\label{subsec:phys-req}

\begin{itemize}
    \item \textbf{REQ\_06\_PHYS}: [Requirement description.]
    \item \textbf{REQ\_07\_PHYS}: [Requirement description.]
\end{itemize}

\subsection{Non-Functional Requirements}
\label{subsec:nonfunc-req}

\begin{itemize}
    \item \textbf{REQ\_08\_USABILITY}: [Requirement description.]
    \item \textbf{REQ\_09\_MODULARITY}: [Requirement description.]
\end{itemize}

\section{Risk Analysis}
\label{sec:risk-analysis}

[Summarise the risk analysis methodology and key findings. Reference
the detailed risk register in Appendix~\ref{app:risk}.]

% ============================================================
%  CHAPTER 3 – LITERATURE STUDY
% ============================================================
\chapter{Literature Study}
\label{ch:literature}

\section{Case Studies}
\label{sec:case-studies}

[Review relevant published work, prior art, and case studies. Each
sub-section can address a different case or research area.]

\subsection{Case Study 1 – [Topic]}

[Discussion with appropriate citations~\cite{ref1}.]

\subsection{Case Study 2 – [Topic]}

[Discussion with appropriate citations~\cite{ref2}.]

\section{Existing Solutions}
\label{sec:existing-solutions}

[Compare and critique existing commercial or academic solutions
relevant to your project.]

\section{Possible Solutions}
\label{sec:possible-solutions}

[Identify two or three candidate approaches. Discuss trade-offs and
justify the selected approach.]

\section{Summary}
\label{sec:lit-summary}

[Brief paragraph drawing together insights from the literature that
informed the design decisions.]

% ============================================================
%  CHAPTER 4 – CONCEPT DESIGN
% ============================================================
\chapter{Concept Design}
\label{ch:concept}

[Present the high-level functional architecture of the solution.
Use functional flow diagrams and operational level diagrams to
describe what the system must do without specifying how.]

\subsection*{Functional Flow Diagram}

\begin{figure}[H]
    \centering
    %\includegraphics[width=0.75\textwidth]{figures/functional_flow.png}
    \caption{Functional flow diagram illustrating the life-cycle.}
    \label{fig:functional-flow}
\end{figure}

[Table~\ref{tab:functional-units} lists the functional units
identified in the concept design.]

\begin{table}[H]
    \centering
    \caption{Operational Functional Units}
    \label{tab:functional-units}
    \begin{tabular}{@{} l l l @{}}
        \toprule
        \textbf{F/U ID} & \textbf{Name} & \textbf{Description} \\
        \midrule
        F/U 1.0 & [Name] & [Description] \\
        F/U 2.0 & [Name] & [Description] \\
        F/U 3.0 & [Name] & [Description] \\
        \bottomrule
    \end{tabular}
\end{table}

% ============================================================
%  CHAPTER 5 – PRELIMINARY DESIGN
% ============================================================
\chapter{Preliminary Design}
\label{ch:preliminary}

[Present the architectural decisions made at the subsystem level.
Describe which technologies, components, or strategies were considered
for each functional unit and justify the selections.]

\section{Preliminary Hardware Decisions}
\label{sec:prelim-hw}

\subsection{[Subsystem 1 – e.g. Processing Unit]}

[Discussion of candidate components. Summarise in a comparison table.]

\begin{table}[H]
    \centering
    \caption{[Subsystem 1] Component Comparison}
    \label{tab:subsystem1-compare}
    \begin{tabular}{@{} l c c c @{}}
        \toprule
        \textbf{Component} & \textbf{Criterion A} & \textbf{Criterion B} & \textbf{Selected} \\
        \midrule
        Option 1 & & & \\
        Option 2 & & & \checkmark \\
        Option 3 & & & \\
        \bottomrule
    \end{tabular}
\end{table}

\subsection{[Subsystem 2]}

[Discussion and selected component with justification~\cite{ref3}.]

\begin{figure}[H]
    \centering
    %\includegraphics[width=0.45\textwidth]{figures/component_photo.png}
    \caption{[Selected component name and description].}
    \label{fig:component}
\end{figure}

% ============================================================
%  CHAPTER 6 – DETAILED DESIGN
% ============================================================
\chapter{Detailed Design}
\label{ch:detailed}

\section{Detailed Hardware Decisions}
\label{sec:detailed-hw}

[Specify exact components, part numbers, electrical parameters, and
interfaces. Include schematics and specification tables.]

\begin{table}[H]
    \centering
    \caption{[Component] Specifications}
    \label{tab:component-specs}
    \begin{tabular}{@{} l l @{}}
        \toprule
        \textbf{Parameter} & \textbf{Value} \\
        \midrule
        Supply voltage      & \SI{5}{\volt} \\
        Operating current   & \SI{500}{\milli\ampere} \\
        Dimensions          & \SI{85}{\milli\metre} $\times$ \SI{56}{\milli\metre} \\
        Mass                & \SI{45}{\gram} \\
        \bottomrule
    \end{tabular}
\end{table}

\section{Detailed Software Decisions}
\label{sec:detailed-sw}

[Describe software architecture, programming languages, libraries,
algorithms, and key design patterns selected.]

\section{CAD Designs}
\label{sec:cad}

[Present mechanical CAD drawings for any custom-fabricated parts.
Include isometric views and critical dimensions.]

\begin{figure}[H]
    \centering
   % \includegraphics[width=0.55\textwidth]{figures/cad_isometric.png}
    \caption{Isometric CAD drawing of [assembly name].}
    \label{fig:cad-isometric}
\end{figure}

% ============================================================
%  CHAPTER 7 – IMPLEMENTATION
% ============================================================
\chapter{Implementation}
\label{ch:implementation}

\section{Electrical Implementation}
\label{sec:elec-impl}

[Describe the electrical assembly process, wiring, PCB fabrication,
and any calibration steps. Include wiring diagrams.]

\begin{figure}[H]
    \centering
    %\includegraphics[width=0.85\textwidth]{figures/wiring_diagram.png}
    \caption{Electrical hardware component layout and wiring diagram.}
    \label{fig:wiring}
\end{figure}

\section{Mechanical Implementation}
\label{sec:mech-impl}

[Describe the physical construction, materials, machining, and
assembly of mechanical components.]

\section{Software Implementation}
\label{sec:sw-impl}

[Describe how the software was implemented. Include a high-level
flow diagram and, where appropriate, annotated code snippets.]

\begin{lstlisting}[language=Python, caption={Example: sensor initialisation routine.}, label=lst:sensor-init]
import sensor_library as sl

def init_sensor(pin: int, baud: int = 9600) -> sl.Sensor:
    """Initialise and calibrate the sensor on the given pin."""
    sensor = sl.Sensor(pin=pin, baud=baud)
    sensor.calibrate()
    return sensor
\end{lstlisting}

\section{Final System}
\label{sec:final-system}

[Present photographs of the completed system and a brief description
of its final form.]

\begin{figure}[H]
    \centering
    \begin{subfigure}[b]{0.45\textwidth}
       % \includegraphics[width=\textwidth]{figures/final_front.png}
        \caption{Front view.}
    \end{subfigure}
    \hfill
    \begin{subfigure}[b]{0.45\textwidth}
        %\includegraphics[width=\textwidth]{figures/final_side.png}
        \caption{Side view.}
    \end{subfigure}
    \caption{Final assembled system.}
    \label{fig:final-system}
\end{figure}

% ============================================================
%  CHAPTER 8 – TESTING
% ============================================================
\chapter{Testing}
\label{ch:testing}

[Each test should trace back to a specific requirement. State the
test procedure, expected outcome, measured result, and pass/fail
verdict.]

\section{Test REQ\_01\_FUNC – [Requirement Name]}
\label{sec:test-req01}

\textbf{Objective:} Verify that \ldots

\textbf{Procedure:}
\begin{enumerate}
    \item Step 1 \ldots
    \item Step 2 \ldots
    \item Record measurement.
\end{enumerate}

\textbf{Results:}

\begin{table}[H]
    \centering
    \caption{Results for Test REQ\_01\_FUNC}
    \label{tab:test-req01}
    \begin{tabular}{@{} l c c c @{}}
        \toprule
        \textbf{Trial} & \textbf{Expected} & \textbf{Measured} & \textbf{Pass/Fail} \\
        \midrule
        1 & & & Pass \\
        2 & & & Pass \\
        3 & & & Pass \\
        \bottomrule
    \end{tabular}
\end{table}

\textbf{Conclusion:} The requirement was [met / not met] because \ldots

\section{Test REQ\_04\_PERF – [Requirement Name]}
\label{sec:test-req04}

[Repeat pattern for each requirement under test.]

% ============================================================
%  CHAPTER 9 – CONCLUSION
% ============================================================
\chapter{Conclusion}
\label{ch:conclusion}

[Summarise what was achieved relative to the objectives stated in
Chapter~\ref{ch:problem}. Discuss limitations and suggest avenues
for future work.]

% ============================================================
%  REFERENCES
% ============================================================
\renewcommand{\bibname}{References}
\begin{thebibliography}{99}

\bibitem{ref1}
A.~Author, B.~Author,
``Title of journal article,''
\textit{Journal Name}, vol.~X, no.~Y, pp.~ZZ--ZZ, Year.
doi: \texttt{10.xxxx/xxxxxxx}.

\bibitem{ref2}
C.~Author,
\textit{Title of Book}.
City: Publisher, Year.

\bibitem{ref3}
D.~Author, E.~Author, and F.~Author,
``Title of conference paper,''
in \textit{Proc. Conference Name}, City, Country, Year, pp.~ZZ--ZZ.

\bibitem{ref4}
Manufacturer Name,
\textit{Product Datasheet: [Part Number]},
[Online]. Available: \url{https://www.example.com/datasheet.pdf}.
[Accessed: Day Month Year].

\end{thebibliography}

% ============================================================
%  APPENDICES
% ============================================================
\begin{appendices}
\renewcommand{\chaptername}{Appendix}

% ─── Appendix A ─────────────────────────────────────────────
\chapter{Generative AI Usage}
\label{app:ai-usage}

\begin{table}[H]
    \centering
    \caption{Generative AI (GA) evidence throughout the report}
    \label{tab:ga-evidence}
    \begin{tabular}{@{} l l l p{5cm} @{}}
        \toprule
        \textbf{Section} & \textbf{Tool} & \textbf{Prompt} & \textbf{How output was used} \\
        \midrule
       % [e.g.~\S\ref{sec:sw-impl}] & ChatGPT & ``\ldots'' & Edited and verified before inclusion. \\
        \bottomrule
    \end{tabular}
\end{table}

% ─── Appendix B ─────────────────────────────────────────────
\chapter{Risk Register}
\label{app:risk}

\begin{table}[H]
    \centering
    \caption{Risk Identification and Management}
    \label{tab:risk-register}
    \begin{tabular}{@{} l p{3cm} l l p{3.5cm} @{}}
        \toprule
        \textbf{ID} & \textbf{Risk} & \textbf{Likelihood} & \textbf{Impact} & \textbf{Mitigation} \\
        \midrule
        R01 & [Description] & Low   & High   & [Mitigation strategy] \\
        R02 & [Description] & Medium & Medium & [Mitigation strategy] \\
        R03 & [Description] & High  & Low    & [Mitigation strategy] \\
        \bottomrule
    \end{tabular}
\end{table}

\begin{table}[H]
    \centering
    \caption{Risk Traceability Matrix BEFORE Mitigation}
    \label{tab:risk-before}
    \begin{tabular}{@{} l ccccc @{}}
        \toprule
         & \multicolumn{5}{c}{\textbf{Impact}} \\
        \textbf{Likelihood} & \textbf{1} & \textbf{2} & \textbf{3} & \textbf{4} & \textbf{5} \\
        \midrule
        5 (Very High) & & & & & \\
        4 (High)      & & & & & \\
        3 (Medium)    & & & & & \\
        2 (Low)       & & & & & \\
        1 (Very Low)  & & & & & \\
        \bottomrule
    \end{tabular}
\end{table}

% ─── Appendix C ─────────────────────────────────────────────
\chapter{Additional Technical Data}
\label{app:technical}

[Include any supplementary data sheets, extended test results,
circuit diagrams, or source-code listings that are too detailed
for the main body but are needed for completeness.]

\end{appendices}

% ============================================================
\end{document}
% ============================================================
